\chapter{\sentil{} in other languages}
\label{ch:bindings}

As mentioned above, \sentil{} is a Rust core wrapped in a thin, stable C interface, and every binding calls that same core. This makes it possible to extend \sentil{} to many languages without reimplementing the engine, and to keep them all in sync.

\section{\sentil{} in Python, C++, Julia, and Java}

The trace is \code{speed = [12, 9, 7, 4, 6]}, the spec is \code{G (speed > 5)}, and the answer is $-1$, driven by the dip to $4$ at $t=3$. Here it is in Rust and in four bindings;

\begin{lstlisting}[language=sentilrust]
// Rust
let mut trace = Trace::new(vec![0.0, 1.0, 2.0, 3.0, 4.0])?;
trace.add_signal("speed", vec![12.0, 9.0, 7.0, 4.0, 6.0])?;
let phi = Formula::parse("G (speed > 5)")?;
phi.robustness(&trace)?;          // -1
phi.robustness_signal(&trace)?;   // [-1.0, -1.0, -1.0, -1.0, 1.0]
\end{lstlisting}

\begin{lstlisting}[language=Python]
# Python
trace = sentil.Trace([0, 1, 2, 3, 4], {"speed": [12., 9., 7., 4., 6.]})
phi = Formula.parse("G (speed > 5)")
phi.robustness(trace)                                      # -1.0
[(v.start, v.end) for v in phi.violations(trace)]          # [(0.0, 3.0)]
\end{lstlisting}

\begin{lstlisting}[language=C++]
// C++
sentil::Trace trace({0, 1, 2, 3, 4}, "speed", {12., 9., 7., 4., 6.});
sentil::Formula phi = sentil::Formula::parse("G (speed > 5)");
phi.robustness(trace);            // -1
\end{lstlisting}

\begin{lstlisting}[language=sentiljulia]
# Julia
phi = formula("G (speed > 5)")
trace = Trace(collect(0.0:1.0:4.0), "speed", [12., 9., 7., 4., 6.])
robustness(phi, trace)            # -1.0
\end{lstlisting}

\begin{lstlisting}[language=bash]
# CLI
sentil check -f 'G (speed > 5)' -t trace.csv
#   verdict     violated
#   robustness  -1.000000   (exit 10 = violation, not error)
\end{lstlisting}

\begin{pitfall}
The streaming verdict carries the same five fields everywhere, \code{resolved}, \code{satisfied}, \code{value}, \code{lower}, \code{upper}, but Rust exposes them as methods (\code{verdict.is\_resolved()}) and the others as fields (\code{verdict.resolved}).
\end{pitfall}

\section{Feature Parity}

Every binding ships a \code{test\_oracle} that loads the shared \code{benchmarks/deterministic} \code{/oracle.json} and reproduces each stored value so you can be sure that the FFI didn't break anything.

\begin{table}[h]
\centering\small
\renewcommand{\arraystretch}{1.3}
\begin{widefig}\begin{tabular}{@{}l l l@{}}
\toprule
\textbf{operation} & \textbf{Python} & \textbf{Rust}\\
\midrule
parse a formula & \code{Formula.parse(s)} & \code{Formula::parse(s)}\\
offline robustness & \code{phi.robustness(tr)} & \code{phi.robustness(\&tr)}\\
stream a sample & \code{mon.update(t, m)} & \code{mon.update(t, \&m)}\\
probabilistic check & \code{phi.check(tr, lift, cfg)} & \code{phi.check(\&tr, \&lift, \&cfg)}\\
synthesize an input & \code{synthesis.synthesize(\dots)} & \code{Synthesizer::solve(\&prob)}\\
\bottomrule
\end{tabular}\end{widefig}
\caption{The same operations under each language's idioms.}
\end{table}

\begin{notebox}
Synthesis is the one place to compare the \emph{objective}, not the answer. Gradient ascent returns any input that reaches the optimum, so Python may find \code{[1, -1, 0]} and Rust \code{[1, 1, 0.41]} for the same problem, both scoring robustness $1.0$. The achieved robustness matches across bindings; the witness input need not.
\end{notebox}

\section{\sentil{} on microcontrollers}

Every binding covers the full engine with the exception of the Arduino target. This is because it cannot run statistical model checking or the GPU path, because a microcontroller has neither the memory for an ensemble nor a GPU. It delivers the streaming deterministic monitor in full though.

\section{FFI Boundary Costs}

Calling across a language boundary is not free, and \sentil{} keeps the cost small by keeping the boundary thin: primitives and opaque handles cross it, the heavy work stays in Rust. Table~\ref{tab:overhead} is the per-sample cost of one streaming update through each binding, on the same nested formula.

\begin{table}[h]
\centering\small
\renewcommand{\arraystretch}{1.25}
\begin{tabular}{@{}l r@{}}
\toprule
\textbf{binding} & \textbf{per update}\\
\midrule
C            & 74 ns\\
C++          & 134 ns \\
Rust (core)  & 112 ns \\
Python       & 521 ns\\
Julia        & 114 ns \\
Java         & 680 ns\\
MATLAB       & 6.49 \textmu s\\
\bottomrule
\end{tabular}
\caption{Per-sample streaming cost by binding, same workload and machine.}
\label{tab:overhead}
\end{table}

Because the core is shared, a correctness fix or a speedup lands for every language at once. There is no per-binding reimplementation to drift out of sync, which is the usual way multi-language tools grow inconsistent.